\documentclass{article}
\input{preamble}
\input{filewise_preamble}
\title{A Sample Paper for Leanization}
\author{Your Name}
\date{}
\begin{document}
\maketitle
\begin{abstract}
This example demonstrates the correspondence between a mathematical paper
and a Lean formalization. Replace it with your own manuscript.
\end{abstract}
\section{Definitions and a first result}
\begin{definition}[\leanlink{def:twice}{Twice}]\label{def:twice}
For a natural number $n$, define $\operatorname{twice}(n)=n+n$.
\end{definition}
\begin{theorem}[\leanlink{thm:twice-eq-add-self}{Expansion}]\label{thm:expansion}
For every natural number $n$, we have $\operatorname{twice}(n)=n+n$.
\end{theorem}
\begin{proof}
This follows immediately from Definition~\ref{def:twice}.
\end{proof}
\end{document}
